% TOP_PROBLEM: {"id": 3700004, "title": "Bounded wandering domains of entire functions", "category_id": 11, "category": {"id": 11, "name": "dynamical_systems", "display_name": "Dynamical Systems", "description": "Problems about long-term behavior of deterministic systems, Hamiltonian dynamics, and periodic orbits.", "slug": "dynamical-systems", "order_index": 11, "created_at": "2026-07-31T15:26:25.671Z"}, "status": "open"}
% FIRST_POSTED: null
% ENTRY_KIND: statement_only
% Imported from: batch12.tex; UnsolvedMath ID 3700004
% Compile twice: lualatex 3700004.tex
\documentclass[11pt]{article}
\usepackage{fontspec}
\setmainfont{Latin Modern Roman}
\usepackage{amsmath,amssymb,mathtools,unicode-math}
\setmathfont{Latin Modern Math}
\usepackage{xurl,hyperref}
\hypersetup{hidelinks,pdftitle={UnsolvedMath Problem 3700004}}
\newcommand{\sref}[2]{\hyperlink{source:#1:#2}{[S#2]}}
\newcommand{\cataloguescope}[1]{\paragraph{Mathematical question.}#1}
\begin{document}
% UnsolvedMath ID 3700004; source catalogue code AMR-036-0004
\setcounter{section}{3700003}
\section{Bounded wandering domains of entire functions}\label{3700004}
{\small\sffamily UnsolvedMath ID 3700004 · Dynamical Systems}
\subsection{Definitions and mathematical statement}

\paragraph{Shared notation.}$\mathbb N=\{1,2,\ldots\}$, and $\mathbb Z,\mathbb Q,\mathbb R,\mathbb C$ denote the integers, rationals, reals and complex numbers. Any other conventions are specified locally.


For a nonlinear entire function $f:\mathbb C\to\mathbb C$, its Fatou set is the open set on which the iterates $\{f^n:n\ge0\}$ form a normal family for locally uniform convergence in the spherical metric. Let $D$ be a connected component of that set. A constant limit function is a value $c\in\widehat{\mathbb C}$ for which some subsequence $f^{n_k}$ tends locally uniformly on $D$ to the constant $c$, with $n_k\to\infty$. Suppose every subsequential limit function on $D$ is constant, and put
\[
\mathcal L(D)=\{c\in\widehat{\mathbb C}: f^{n_k}|_D\longrightarrow c\text{ locally uniformly for some }n_k\to\infty\}.
\]
A component is wandering if the Fatou components containing $f^n(D)$ are pairwise distinct. Boundedness of $\mathcal L(D)$ means $\mathcal L(D)\subset\mathbb C$ and $\sup_{c\in\mathcal L(D)}|c|<\infty$, so infinity is excluded.
\paragraph{Statement recorded in the supplied source.}
Can $\mathcal L(D)$ be both infinite and bounded? Equivalently in the author's formulation, can there be a nonempty open subdomain $U\subset D$ of a wandering component whose forward images remain in a bounded subset of $\mathbb C$?

The September 2026 paper of Drach, Pardo-Simón and Učakar states a negative answer in Theorem B: for every point $z$ of a wandering component of a transcendental meromorphic function, $\limsup_{n\to\infty}|f^n(z)|=\infty$. Its historical discussion explicitly identifies the infinite bounded constant-limit-set question. \sref{3700004}{2} \sref{3700004}{3}
\subsection{Short English statement}
Can an entire function have a Fatou component with infinitely many constant subsequential limits, all in a bounded part of the plane? The cited 2026 theorem reports that bounded wandering orbits do not occur.
\subsection{Sources}
\begin{description}
\item[\hypertarget{source:3700004:1}{S1}] ULAM.AI and the UnsolvedMath Contributors, UnsolvedMath: A Curated Collection of Open Mathematics Problems, record 3700004 (AMR-036-0004). CC BY 4.0. \url{https://huggingface.co/datasets/ulamai/UnsolvedMath}.
\item[\hypertarget{source:3700004:2}{S2}] Ben Bielefeld and Mikhail Lyubich (editors); Alexandre Eremenko and Mikhail Lyubich (problem authors), Problems in Holomorphic Dynamics (1992), Eremenko–Lyubich, Wandering Domains for Holomorphic Maps, Question 1, pp.39–40. \url{https://arxiv.org/abs/math/9205209}.
\item[\hypertarget{source:3700004:3}{S3}] Kostiantyn Drach, Leticia Pardo-Simón and Beno Učakar, Bounded-orbit wandering domains do not exist (2026), Theorems A,B; Remark 1.1; §1.1, pp.1–3. \url{https://arxiv.org/abs/2609.30103}.
\end{description}
\label{end:3700004}
\end{document}
